\documentclass[11pt]{article}
\usepackage{amsmath,amsfonts}
\usepackage{graphicx}
%\usepackage{xcolor}
%\definecolor{gray}{rgb}{.75,.75,.75}
%\usepackage[landscape]{geometry}
%\usepackage{hyperref}
%\usepackage[bookmarks=true, pdfborder={0 0 .5}, urlbordercolor=gray]{hyperref}

\renewcommand{\thefootnote}{\fnsymbol{footnote}}

\addtolength{\topmargin}{-.5in}
\addtolength{\textheight}{.9in}
\addtolength{\oddsidemargin}{-.3in}
\addtolength{\textwidth}{.6in}
\pagestyle{empty}

\newcommand{\ds}{\displaystyle}
\newcommand{\ts}{\textstyle}

\newcommand{\Z}{\mathbb{Z}}
\newcommand{\R}{\mathbb{R}}
\newcommand{\C}{\mathbb{C}}
\newcommand{\EE}{\mathcal{E}}
\newcommand{\tZ}{\tilde{\Z}}

\newcommand{\sech}{\mathrm{sech}}
\newcommand{\csch}{\mathrm{csch}}


\newcommand{\Sym}{\mathrm{Sym}}

\renewcommand{\circle}[1]{{\bigcirc\hspace{-.75em}#1}}

\begin{document}

\footnote[0]{Notes by Peter Magyar \ {\tt magyar@math.msu.edu}}
\vspace{-1em}

\centerline{\bf Math 133\hspace{1.5em} \hfill{\Large\bf
 Taylor Series}\hfill Stewart \S11.10}
\vspace{1em}

\noindent 
{\bf Series representation of a function.} A series writes a given {\it complicated} quantity as an infinite sum of {\it simple} terms.
To approximate the quantity, we take only the first few terms of the series,
dropping the later terms which give smaller and smaller corrections. 
In this section, we finally develop the tool that lets us do this in most cases: a way to write any reasonable function $f(x)$ as an explicit power series, a kind of infinte polynomial. 
This will allow us to compute outputs of the function 
by plugging values of $x$ into the series.

Our functions must behave decently near the center point of the desired power series.
We say $f(x)$ is {\it analytic} at $x=a$ if it is possible to write $f(x)=\sum_{n=0}^\infty c_n(x{-}a)^n$ for some coefficients $c_n$, with positive radius of convergence.
In practice, any formula involving standard functions and operations 
defines an analytic function, provided the formula gives real number values in a
small interval around $x=a$. For example $\frac{1}{x{-}a}$ is {\it not} analytic at $x=a$, 
because it gives $\pm\infty$ at $x=a$; and $\sqrt{x{-}a}$ is {\it not} analytic at $x=a$ because
for $x$ slightly smaller than $a$, it gives the square root of a negative number.\footnote{
The function $\sqrt[3]{x{-}a\,}$ is also not analytic near $x=a$,
ever though it gives real number values.
The problem is it has a vertical tangent at $x=a$,
so it is not differentiable. Also see $e^{-1/x^2}$ in \S11.11.
} 
\begin{quote}
{\it Taylor Series Theorem:} Let $f(x)$ be a function which is analytic at $x=a$.
Then we can write $f(x)$ as the following power series,  called the {\it Taylor series} of $f(x)$
at $x=a$:
$$
f(x) \ =\ f(a) + f'(a)\,(x{-}a) + \frac{f''(a)}{2!}\,(x{-}a)^2 + \frac{f'''(a)}{3!}\,(x{-}a)^3 + \frac{f''''(a)}{4!}\,(x{-}a)^4 + \cdots\,,
$$
valid for $x$ within a radius of convergence $|x{-}a|<R$ with $R>0$, or convergent for all $x$.

If we write the $n$th derivative of $f(x)$ as $f^{(n)}(x)$, this becomes: 
$$
f(x) \ =\  \sum_{n=0}^\infty c_n(x{-}a)^n
\qquad
\text{with coefficients}
\qquad
c_n=\frac{f^{(n)}(a)}{n!}\,.
$$
\end{quote}
%
{\sc warning:} The coefficients are {\it constants} with no $x$, 
so $c_1 =  f'(a)$, {\sc not} $ f'(x)$.
\\[.5em]
{\it Proof.} By hypothesis $f(x)$ is analytic, so $f(x)=\sum_{n=0}^\infty c_n(x{-}a)^n$ 
for {\it some} $c_n$; we will derive the desired formula for these coefficients. Since $f(a) = \sum_{n=0}^\infty c_n(a{-}a)^n=
c_0 + c_1(0) + c_2(0^2)+\cdots$, we get $c_0=f(a)$. 
Next, by the Theorem in \S11.9, we have (dropping the zero term)
$f'(x) = \sum_{n=1}^\infty nc_n (x{-}a)^{n-1}$, 
so $f'(a) = c_1 + 2c_2(0) + 3c_3(0^2) + \cdots$, and $c_1 = f'(a)$. 
Next, $f''(x) = \sum_{n=2}^\infty n(n{-}1)c_n (x{-}a)^{n-2}$,
 so $f''(a) = (2)(1)c_2$ and $c_2=\frac12 f''(a)$.  Continuing, we get:
$$
f^{(N)}(x) \ =\ \sum_{n=N}^\infty n(n{-}1)\cdots(n{-}N{+}1)\, c_n\,(x{-}a)^{n-N}\,.
$$
The terms for $n= 0,1,\ldots, N{-}1$ are all zero, as is evident from the factors $n(n{-}1)\cdots(n{-}N+1)$,
so the first non-zero term is for $n=N$. Plugging in $x=a$ gives:
$f^{(N)}(a) \ = \ N(N{-}1)\cdots(1)c_N$, and $c_N=\frac{1}{N!} f^{(N)}(a)$ as desired. Q.E.D.

Once we have a power series for $f(x)$ with known coefficients $c_n = \frac{f^{(n)}(a)}{n!}$, 
we can approximate $f(x)$ by taking a finite partial sum of the series up to some cutoff term $N$. 
This partial sum is called a {\it Taylor polynomial}, denoted $T_N(x)$:
$$
f(x)\ \approx\ T_N(x) \ =\ \sum_{n=0}^N c_n(x{-}a)^n = f(a) + f'(a)(x{-}a)  + \cdots + \frac{f^{(N)}(a)}{N!}(x{-}a)^N\,.
$$
Note that $T_1(x) = f(a)+f'(a)(x-a)$ is just the linear approximation near $x=a$, whose graph is the tangent line (Calculus I \S2.9).
We can improve this approximation of $f(x)$ in two ways:
\begin{itemize}
\item Take more terms, increasing $N$.
\item Take  the center $a$ close to $x$, giving small $(x{-}a)$ and tiny $(x{-}a)^n$.
\end{itemize}
%
A Taylor series centered at $a=0$ is specially named a {\it Maclaurin series}.
\\[1em]
{\bf Example: sine function.} To find Taylor series for a function $f(x)$, we must determine
$f^{(n)}(a)$. This is easiest for a function which satisfies a simple differential equation
relating the derivatives to the original function. For example, $f(x)=\sin(x)$ satisfies 
$f''(x) = -f(x)$, so coefficients of the Maclaurin series (center $a=0$) are:
$$\begin{array}{|c||c|c|c|c|c|c|c|c|c|}
\hline &&&&&&&& \\[-.9em] 
n & 0 & 1 & 2  & 3 & 4 & 5 & 6 & 7 
\\[.2em] \hline &&&&&&&& \\[-.9em]
f^{(n)}(x) & \,\sin(x)\, &\,\cos(x)\, & \!-\sin(x)\! & \!-\cos(x)\! & \,\sin(x)\, &\,\cos(x)\, &\! -\sin(x)\!& \!-\cos(x) \!
\\[.2em] \hline &&&&&&&& \\[-.9em] 
f^{(n)}(0) & 0 &1 & 0 & -1 & 0 &1 & 0 & -1 
\\[.2em] \hline &&&&&&&&\\[-.9em] 
\!c_n\,{=}\,\frac{f^{(n)}(0)}{n!}\! & 0 & 1 & 0 & -\frac1{3!} & 0 &\frac1{5!} & 0 &-\frac1{7!}
\\[.3em] \hline
\end{array}$$
That is:
$$
\sin(x) \ =\  x - \frac{x^3}{3!} +\frac{x^5}{5!}- \frac{x^7}{7!}+\cdots 
\ =\ \sum_{n=0}^\infty (-1)^n \frac{x^{2n+1}}{(2n{+}1)!}\,. 
$$
To find the domain of convergence, we apply the Ratio Test (11.6/I):
$$
L 
\ = \ \lim_{n\to\infty}\left|\frac{a_{n+1}}{a_n}\right|
\ = \ \lim_{n\to\infty}\left|\frac{x^{2(n+1)+1}}{(2(n{+}1){+}1)!}\ \left/ \frac{x^{2n+1}}{(2n{+}1)!}\right.\right|
$$ 
$$
\ = \ \lim_{n\to\infty}\frac{|x|^{2n+3}}{|x|^{2n+1}}\cdot\frac{(2n{+}1)!}{(2n{+}3)!}
\ = \ \lim_{n\to\infty}\frac{|x|^2}{(2n{+}2)(2n{+}3)}
\ =\ 0
$$ 
for any fixed $x\neq 0$. Since $L=0<1$ regardless of $x$, the series converges for all $x$.

This formula for $\sin(x)$ astonishes because the right side is a simple algebraic series  having no apparent relation to trigonometry. We can try to understand and check the series by graphically comparing 
$\sin(x)$ with its Taylor polynomial approximations:
\\[1em]
\centerline{
\includegraphics[width=4.5in]
{11-10-Sin-Taylor.png}
}
\vspace{-1.5em}

\begin{itemize}
\item The Taylor polynomial $T_1(x)=x$ (in red) is just the linear approximation or tangent line 
of $y=\sin(x)$ at the center point $x=0$.
The curve and line are close (to within a couple of decimal places) near the point of tangency
and up to about $|x|\leq 0.5$. Once they veer apart, the approximation is useless.

\item The next Taylor polynomial $T_3(x) = x-\frac{x^3}{3!}=x-\frac16 x^3$ (in green) matches $y=\sin(x)$ in its first three derivatives at $x=0$, and stays close to the original curve up to about $|x|\leq 1.5$\,.

\item The next $T_5(x)= x-\frac{x^3}{3!}+\frac{x^5}{5!}=x-\frac16 x^3+\frac{1}{120}x^5$ is even closer to $f(x)$ for even larger $x$.
Taking enough terms in the Taylor series will give a good approximation for {\it any} $x$, since
the series converges everywhere.
\end{itemize}


\noindent
%A given approximation $f(x)\approx T_N(x)$ is more accurate for smaller $|x|$. In fact, we need not consider large $x$, since all values of $\sin(x)$ can be reduced to $x\in[0,\frac\pi2]$.
%\\[.5em]
{\sc problem:} Compute $\sin(10^\circ)$. A geometric method would be to construct a right triangle with a $10^\circ$ angle, and measure the opposite side divided by the hypotenuse; but this would 
only produce a couple of decimal places of accuracy.
Of course,~a calculator can produce many decimal places, but how does it know?
Taylor series!

As always when doing calculus on trig functions, we must first convert to radians (see end of \S2.5):
$10^\circ = \frac{2\pi}{360}( 10) = \frac{\pi}{18}$. Here $|x|=\tfrac{\pi}{18}\approx\frac16$ is small, so the Maclaurin series centered at $0$ should converge quickly, giving very accurate approximations:
$$
\sin(\tfrac{\pi}{18}) \ \approx\ T_3(\tfrac{\pi}{18}) \ =\  \tfrac{\pi}{18} - \tfrac16(\tfrac{\pi}{18})^3 \ \approx\ \underline{0.17364}68\,.
$$
It turns out this is correct to 5 decimal places (underlined), 
using only two non-zero terms of the Taylor series and a good estimate for $\pi$.
We could verify this by taking more terms and seeing that these 5 digits do not change,
or by applying the Lagrange Remainder estimates in \S11.11.
\newpage

\noindent
{\bf Example: square roots.} Compute $\sqrt{2}$ to 5 decimal places.\footnote{
We saw a faster algorithm for this (but not for functions like $\sin$) in Calculus I \S3.8: 
Newton's Method finds approximate solutions to equations like $g(x)=x^2-2=0$ 
by repeatedly solving a linear approximation of $g(x)=0$;
this improves approximate solution $x_n$ to $x_{n+1}=x_n-g(x_n)/g'(x_n)
=x_n-(x_n^2{-}2)/(2x_n)$.
Starting with $x_0=3/2=1.5$ gives $x=\sqrt 2$ accurate to 5 decimal places
after just $n=2$ iterations, while the Taylor series requires $n=4$.
Newton's Method doubles the number of accurate places each time, 
while each term of the series adds a constant number of places (\S11.11).
} 
First, we must consider $\sqrt 2$ to be an output of the function $f(x)=\sqrt x$ at $x=2$.
Next, we must choose the center $a$ for its Taylor series. 
\begin{itemize}
\item $a=0$ does not give a series because $\sqrt x$ is not analytic at $x=0$.
Indeed, if there were a convergent Taylor series $\sqrt x= c_0 + c_1x+ c_2x^2+\cdots$, 
we could plug in $x=-0.1$ to get:
$\sqrt{-0.1}=c_0 + c_1(-0.1)+ c_2(-0.1)^2+\cdots$, a real value
for the square root of a negative number! 

\item $a=1$ is too far from $x=2$: it turns out $|x{-}a| = |2{-}1|= 1$ is
beyond the radius of convergence of the Taylor series.

\item $a=2$ is useless, since writing the Taylor series requires us to know $f^{(n)}(2)$, including $f(2)=\sqrt 2$, the same number 
we are trying to compute.

\item A useful choice of $a$ requires: $a>0$  so that the Taylor series exists; 
$a$ is close to $x=2$, making $|x{-}a|$ small so the series converges quickly;
and $f(a)=\sqrt a$ is easy to compute  so we can find the coefficients.
A value satisfying all three conditions is:
$a=\tfrac94\,.$

\end{itemize}

Now we have:
$$\begin{array}{|c||c|c|c|c|c|}
\hline &&&&& \\[-.9em] 
n & 0 & 1 & 2  & 3 & 4 
\\[.2em] \hline &&&&& \\[-.9em]
f^{(n)}(x) & x^{1/2} & \frac12\,x^{-1/2} 
&-\frac{1\cdot 1}{2\cdot 2}\,x^{-3/2} 
& \frac{1\cdot 1\cdot 3}{2\cdot 2\cdot 2}\,x^{-5/2} 
& -\frac{1\cdot 1\cdot 3\cdot 5}{2\cdot 2\cdot 2\cdot 2}\,x^{-7/2} 
\\[.2em] \hline &&&&& \\[-.9em] 
f^{(n)}(\frac94) & \frac{3}{2} & \frac{1}{3} & -\frac{2}{27} & \frac{4}{81} & -\frac{40}{729}
\\[.3em] \hline &&&&&\\[-1em] 
\!c_n\,{=}\,\frac{f^{(n)}(\frac94)}{n!}\! & \frac32 &\frac13 & -\frac1{27} & \frac2{243} & -\frac5{2187}  
\\[.5em] \hline
\end{array}$$
Hence:
$$
\sqrt{x} \ =\ 
\tfrac32 + \tfrac13(x{-}\tfrac94)
-\tfrac1{27}(x{-}\tfrac94)^2 + \tfrac{2}{243}(x{-}\tfrac94)^3
- \tfrac{5}{2187}(x{-}\tfrac94)^4+\cdots
$$
$$
\ =\ \tfrac 32 + \tfrac13(x{-}\tfrac94)+ \sum_{n=2}^\infty (-1)^{n-1}\tfrac{(2n-3)!!}{n!}\,\tfrac{2^{n-1}}{3^{2n-1}}\,(x{-}\tfrac94)^n\,,
$$
where we use the odd-factorial notation $(2n{-}3)!! = (1)(3)(5)\cdots
(2n{-}3)$.
For $x=2$, we have $x{-}\tfrac94 = -\frac14$, so:
$$\begin{array}{rcl}
\sqrt 2 &=& \tfrac32 + \tfrac13(-\tfrac14)
-\tfrac1{27}(-\tfrac14)^2 + \tfrac{2}{243}(-\tfrac14)^3
- \tfrac{5}{2187}(-\tfrac14)^4+\cdots
\\[.5em]
&\approx& 
\tfrac32 - \tfrac13\,\tfrac14
-\tfrac1{27}\,\tfrac1{4^2} - \tfrac{2}{243}\,\tfrac1{4^3}
- \tfrac{5}{2187}\tfrac1{4^4}
\ \approx\ \underline{1.41421}43\,,
\end{array}$$
which is correct to 5 decimal places (underlined).
\newpage

\noindent
{\bf Common Taylor series}
\begin{itemize}


\item $\ds\frac{1}{1{-}x} \ =\ \sum_{n=0}^\infty x^n$ \ for $|x|<1$ \ ({\it Geometric Series}).

\item $\ds\ln(1{+}x) \ =\ 
\sum_{n=1}^\infty (-1)^{n-1} \frac{x^n}{n}$ \ for $|x|<1$.

\item $(1{+}x)^p \ =\  \ds\sum_{n=0}^\infty
\frac{p(p{-}1)\cdots (p{-}n{+}1)}{n!}\, x^n$ \  for $|x|<1$
\ ({\it Binomial Series}).

 \item $\ds\exp(x)\ =\ \sum_{n=0}^\infty \frac{x^n}{n!}$ \ for all $x$.  

\item $\ds\sin(x)\ =\ \sum_{n=0}^\infty (-1)^n\frac{x^{2n+1}}{(2n{+}1)!}$ \ for all $x$.

\item $\ds\cos(x)\ =\ \sum_{n=0}^\infty (-1)^n\frac{x^{2n}}{(2n)!}$ \ for all $x$.

\end{itemize}
\vspace{1em}

\noindent
{\bf Euler's Basel Formula.}  
$$ \sum_{n=1}^\infty \frac1{n^2} \ =\  \frac{\pi^2}6.$$
%
To see this, we introduce {\it infinite products}, which represent a function $f(x)$ by
factors of the form $(1-\frac xb)$ where $x=b$ is zero of $f(x)$, a value where $f(b) = 0$.
If $b_1, b_2,\ldots $ is the sequence of all zeroes of $f(x)$, then we should have a 
{\it Weierstrass product}:
$$
f(x) \ = \ c\,\prod_{n=1}^\infty \left(1-\frac x{\ b_n}\right),
$$
where $c=f(0)$, since both sides have the same zeroes and the same value at $x=0$. 

Now apply this to the function $f(x) = \frac{\sin(\pi x)}{\pi x}$, which has Taylor series:
$$
f(x) \ = \ \frac{\sin(\pi x)}{\pi x} 
\ = \ 
\frac{1}{\pi x}\left(\pi x -\frac{(\pi x)^3}{3!} + \frac{(\pi x)^5}{5!}+ \ldots \right)
\ = \ 
1 -\frac{\pi^2}{3!} x^2 + \frac{\pi^4}{5!}x^4 + \ldots   .      
$$
The zeroes are $n = \pm 1, \pm 2, \pm 3,\ldots$, and 
$f(0) = \lim\limits_{x\to0}\frac{\sin(\pi x)}{\pi x}= 1$, so the product is:
$$
f(x) 
\ = \ 
\prod_{n=1}^\infty \left(1-\frac x{n}\right)\left(1+\frac x{n}\right)
\ = \ 
\prod_{n=1}^\infty \left(1-\frac{x^2}{n^2}\right).
$$
Now multiply out the infinite product (taking the limit of the finite products):
$$
f(x) \ = \ 1 -\frac{x^2}{1^2} - \frac{x^2}{2^2} - \frac{x^2}{3^2} - \cdots
+ \frac{x^4}{1^2 2^2} + \frac{x^4}{1^2 3^2} + \cdots
$$
Collecting terms and comparing with the Taylor series:
$$
f(x) \ = \ 1 - \left(\sum_{n=1}^\infty \frac1{n^2}\right)\!x^2 +
\left(\sum_{n=1}^\infty \frac{d(n)}{n^2}\right)\!x^4 + \cdots
\ = \ 1 -\frac{\pi^2}{3!} x^2 + \frac{\pi^4}{5!}x^4 + \cdots .
$$
where $d(n)$ is the number of factors of $n$.
Equating the $x^2$ coefficients gives Euler's formula: $\sum_{i=1}^\infty \frac1{n^2} = \frac{\pi^2}6$\,.
\\[1em]
{\bf Extra Topic: Irrationality of $e$.}  In \S11.2, 
we saw that repeating infinite decimals represent rational numbers (fractions),
and every fraction can be written by long division as a repeating decimal.
Any infinite decimal which does not repeat cannot be written as a fraction: 
it is {\it irrational}. However, it is diffficult to prove that any given interesting number
such as $\pi$ or $\sqrt 2$ is irrational. 

We can use series to prove the irrationality of the constant
$e =2.7182818284590\cdots$.
To prove the negative proposition that $e$ is not equal to any possible fraction $a/b$,
we use the {\it method of contradiction}: 
that is, we assume that there {\it were} some fraction with $e=a/b$, 
and use this to deduce an impossible conclusion, which will show 
that the original assumption  $e=a/b$ was also impossible.

Thus, using the Taylor series definition for $e$, we assume the possibility, 
for some whole numbers $a,b$, of the equation:
$$
1 + \frac{1}{1!} + \frac{1}{2!} +\frac{1}{3!}+ \cdots \ \stackrel{\text{def}}{=}\ e \ =\ \frac ab\,.
$$
It is easy to show $2<e<3$, so $e$ is not a whole number, and would have denominator $b>1$.
Consider the $b^{\mathrm{th}}$ order Taylor approximation,
with error or remainder $R_b$:
$$
1 + \frac{1}{1!} + \frac{1}{2!} +\frac{1}{3!}+ \cdots + \frac{1}{b!}+ R_b \ = \ e,
\qquad
R_b= \sum_{n= b+1}^\infty \frac{1}{n!}.
$$
We multiply by $b!$ to clear denominators up to the $\frac 1{b!}$ term:
$$
b! + \frac{b!}{1!} + \frac{b!}{2!} + \cdots +\frac{b!}{b!}+b!R_b\ \ =\ b!\, e \ =\ b!\,\frac ab \ =\ (b{-}1)!\, a\,.
$$
The terms $b!\,, \frac{b!}{1!}, \frac{b!}{2!}, \ldots ,\frac{b!}{b!}$ on the left
are whole numbers, and $(b{-}1)!\,a$ on the right is a whole number, 
so the remainder $b!R_b$ must also be a whole number.
But it must also be very small,
as we can see from a simple geometric series comparison:\footnote{
Here we do not need the powerful Lagrange remainder formula from \S11.11.
}
$$\begin{array}{rcl}
b!R_b
\ = \  \sum\limits_{n= b+1}^\infty \frac{b!}{n!}
&=&
\frac{1}{b{+}1} + \frac{1}{(b{+}1)(b{+}2)} + \frac{1}{(b{+}1)(b{+}2)(b{+}3)}+\cdots
\\[1em]
&<&
\frac{1}{b+1} + \frac{1}{(b{+}1)^2} + \frac{1}{(b{+}1)^3}+\cdots
\\[1em]
& = & \frac{1}{b{+}1}\frac{1}{1 -(\frac1{b+1})}\ =\ \frac{1}{b{+}1}\frac{b+1}{b+1-1}
\ =\ \frac1b \ < \ 1.
\end{array}$$
Thus, we have constructed a positive number $b!R_b$ which is both a whole number 
and less than 1, which is impossible. 
Therefore the original assumption $e=a/b$ must also be impossible.




\end{document}

%%%%%%%%%%%  Picture  %%%%%%%%%%
\\[0em]
\centerline{
%\includegraphics[width=2in]
{1-8-Discont-iv.png}
\qquad \qquad 
%\includegraphics[width=2in]
{1-8-Discont-v.png} 
}
\vspace{0em}


















